\section{Posterior energy and output ceilings}\label{sec:posterior}
\begin{lemma}[Strong-prior ceiling]\label{lem:strong-ceiling}
Fix a sample, a finite width $N$, depth $L\ge3$, and positive tolerance
$\rho$. Suppose it has an exact interpolant of energy at most $K$.
Then the posterior defined by~\eqref{eq:scales}--\eqref{eq:training}
satisfies
\[
 \Pi_{S,\sigma}\{\kappa_N>K+1\}\longrightarrow0
 \qquad(\sigma\downarrow0).
\]
If $K=2+(L+14)T^{1/(L-1)}$, it consequently assigns probability
tending to one to functions of Lipschitz constant at most
\begin{equation}\label{eq:Lip-ceiling}
 \Lambda=\frac{e^{17/2}}{\sqrt d}T^{3/4}.
\end{equation}
\end{lemma}
\begin{proof}
Continuity gives a parameter ball around the witness that still fits
the sample within $\rho$ and whose energies are at most $K+1/4$.
In the fixed finite parameter dimension its Gaussian probability has
lower exponential rate at least $-(K+1/4)/2$ at speed $\sigma^{-2}$.
The Gaussian radial tail outside energy $K+1$ has upper rate at most
$-(K+1)/2$. These bounds follow directly by integrating the density
over the ball and by the Gaussian squared-norm tail, respectively.
Their quotient tends to zero.

For an energy at most $K+1$, arithmetic--geometric mean and the
one-Lipschitz tanh activation give
\[
 \Lip(z_w)^2\le
 \norm{W^1}_F^2\cdots\norm{W^{L-1}}_F^2\norm v_2^2
 \le\frac1d\left(\frac{K+1}{L}\right)^L
 \le\frac{e^{17}}dT^{3/2},
\]
using~\eqref{eq:product-bound}. The constant input coordinate does
not change the Lipschitz bound with respect to the actual input.
\end{proof}

At finite exponents above one we need a ceiling for a nearby function,
rather than for the full parameter energy. We give its proof explicitly.

\begin{lemma}[Finite-scale approximation ceiling]\label{lem:finite-ceiling}
Suppose that for a sample on the Euclidean unit ball and every
$N\ge M$ there is an exact interpolant of energy at most
$K=2+(L+14)T^{1/(L-1)}$, where $L\ge3$ and $T\ge1$ are independent
of $N$. Fix $\zeta>0$, $0<\beta<1$, and $0<h<\rho$.
Define
\[
 A=\left(\frac{K+2}{L-1}\right)^{(L-1)/2},\qquad
 b^2=\max\{2(d+1)\log2+4\log8,\ 16\log9+8\log(8L)\},
\]
\[
 \eta=\min\left\{1,\frac1{4\sqrt{L(K+1)}},
                         \frac h{\sqrt2 LA}\right\}.
\]
For every
\begin{equation}\label{eq:width-threshold}
 N\ge\left\lceil\max\left\{M,2,(b/\eta)^{2/\zeta},
                         2\log(2/\beta)\right\}\right\rceil,
 \qquad 0<\sigma^2\le N^{-(1+\zeta)},
\end{equation}
the posterior assigns mass at least $1-\beta$ to outputs uniformly
within $h$ on the unit ball of a function with energy at most $K+1$
and Lipschitz constant at most~\eqref{eq:Lip-ceiling}.
The width threshold has a polynomial upper bound in
$d,M,T,L,h^{-1},\log(1/\beta)$ for fixed $\zeta$.
\end{lemma}
\begin{proof}
The energy-coordinate matrices are
\[
 A_1=\sqrt{d/N}W^1,\qquad A_j=W^j\ (2\le j<L),\qquad
 A_L=\sqrt N v^T.
\]
Their entries are independent $\mathcal N(0,\sigma^2)$ and their
total squared Euclidean norm is $\kappa_N$.
After normalizing the augmented input by $\sqrt d$ and hidden states
by $\sqrt N$, the first activation map is
$u\mapsto N^{-1/2}\tanh(\sqrt N A_1u)$, and the other hidden
activation maps have the same form. They are one-Lipschitz in their
linear argument and zero-preserving. The normalized input norm on
the unit ball is at most $\sqrt2$.

Suppose a residual network $U$ has energy at most $K+1$, and a
perturbation $B$ has operator norm at most $\eta$ in each normalized
block. Telescoping one layer at a time and applying
arithmetic--geometric mean to the other $L-1$ layer norms gives
\begin{equation}\label{eq:perturbation}
 \sup_{\norm x_2\le1}|z_{U+B}(x)-z_U(x)|
 \le\sqrt2 L\eta
 \left[\frac{K+1+2\eta\sqrt{L(K+1)}+L\eta^2}{L-1}
 \right]^{(L-1)/2}.
\end{equation}
Indeed, writing the residual block norms as $a_j$, each telescoping
term is at most $\sqrt2\eta$ times a product of $L-1$ numbers
$a_j+\eta$. Their squared sum is bounded by the numerator, since
$\sum a_j^2\le K+1$ and $\sum a_j\le\sqrt{L(K+1)}$.
The choice of $\eta$ makes the cross term at most $1/2$ and
$L\eta^2\le1/[16(K+1)]$. Thus that numerator is at most $K+2$
and~\eqref{eq:perturbation} is at most $\sqrt2L\eta A\le h$.

Put $t^2=N\sigma^2$. Let
\[
 a_d^2=2(d+1)\log2+4\log8,\qquad
 b_L^2=16\log9+8\log(8L).
\]
In energy coordinates let $\mathcal T$ be the set where first and
last block Frobenius norms are at most $ta_d$, and internal block
operator norms are at most $tb_L$. This set is nonempty, closed,
convex and centrally symmetric, and its Gaussian probability is
at least $1/2$. To verify the last assertion, the first squared norm
has law $(t^2/N)\chi^2_{N(d+1)}$; the last uses $N$ degrees of
freedom. The exponential moment at parameter $1/4$ bounds either
failure probability by $1/8$. A $1/4$-net of the unit sphere has at
most $9^N$ points. Applying the scalar Gaussian tail on two such
nets gives, for each internal matrix $tG/\sqrt N$,
\[
 \Prb\{\norm{tG/\sqrt N}_{\rm op}>tb_L\}
 \le2\exp(2N\log9-Nb_L^2/8)\le1/(4L).
\]
The union bound proves the required mass. Every normalized block
in $\mathcal T$ has operator norm at most $tb$.

Condition~\eqref{eq:width-threshold} gives $t\le N^{-\zeta/2}$,
so $tb\le\eta$. If $w_S$ is the exact-fitting witness, then
$w_S+\mathcal T$ fits the sample within $h<\rho$ by the perturbation
bound. The Gaussian shift identity and central symmetry give
\[
 \mu(w_S+\mathcal T)
 =e^{-\kappa_N(w_S)/(2\sigma^2)}
   \int_{\mathcal T}e^{-\langle w_S,B\rangle_\kappa/\sigma^2}\,d\mu(B)
 \ge\tfrac12e^{-K/(2\sigma^2)}.
\]
Pairing $B$ with $-B$ in the integral proves the inequality.

Gaussian isoperimetry~\cite{borell,ledoux}, applied to a set of
Gaussian mass at least $1/2$, gives
\[
 \mu\{\operatorname{dist}_\kappa(w,\mathcal T)^2>K+1\}
 \le e^{-(K+1)/(2\sigma^2)}.
\]
Dividing by the preceding lower bound for the conditioning event
shows that the posterior probability of this event is at most
$2e^{-1/(2\sigma^2)}\le\beta$. Off this event, let $B$ be the metric
projection of $w$ onto $\mathcal T$ and set $U=w-B$.
Then $\kappa_N(U)\le K+1$, and~\eqref{eq:perturbation} gives
$\sup|z_w-z_U|\le h$. The product-norm bound gives the stated
Lipschitz ceiling. No claim that $w$ itself has low energy is needed.

Finally
\[
 K+2\le(L+18)T^{1/(L-1)},\qquad
 A^2\le\left(\frac{L+18}{L-1}\right)^{L-1}T\le e^{19}T.
\]
Thus every inverse radius and threshold in the lemma has the claimed
polynomial upper bound for fixed $\zeta$.
\end{proof}
